\chapter{Materials and Methods}

\section{Data files}
The Flora rests on four tab-separated data files in the \texttt{data}
folder of this project. They are the single source for every count,
date and locality printed in this volume.

\begin{description}[style=nextline,leftmargin=2em]
	\item[\texttt{data/species.tsv}] The checklist: one row for each of
	the 18 species, with its family, author citation, status (native or
	cultivated), the vegetation types in which it occurs, and its local
	and trade names.
	\item[\texttt{data/localities.tsv}] The 12 collecting localities, with
	grid square, altitude and vegetation type. It matches the Gazetteer
	in Appendix~\ref{app:gazetteer} row for row.
	\item[\texttt{data/specimens.tsv}] The 90 specimens on which the
	treatments rest, with collector, collector's number, date, locality,
	altitude, herbarium and phenological state.
	\item[\texttt{data/phenology.tsv}] For each species and month, whether
	the species is recorded in flower (fl), in fruit (fr), both, or
	neither.
\end{description}

\section{Validation of the specimen records}
Every specimen record was checked against six rules before it was
accepted.
\begin{enumerate}[label=\arabic*.]
	\item The altitude equals the altitude of the locality in
	\texttt{localities.tsv}.
	\item According to \texttt{species.tsv}, the species occurs in the
	vegetation type of the locality.
	\item The date falls within the collector's working period: Harding
	1934--1938, Olaniyan 1961--1967, Alade 2019--2025.
	\item The herbarium is the one that holds that collector's material:
	Harding in \textsc{okf}, Olaniyan in \textsc{svn}, Alade in
	\textsc{auh}.
	\item A specimen in flower or fruit was collected in a month in which
	\texttt{phenology.tsv} records that state; a sterile specimen was
	collected in a month with neither.
	\item Each collector's numbers increase with date and are never
	repeated.
\end{enumerate}

\section{Summary of the specimen data}
Table~\ref{tab:family-summary} gives the number of species and
specimens in each family, and Table~\ref{tab:collector-summary} gives
the number of specimens from each collector.

\begin{table}[htbp]
	\centering
	\caption{Species and specimens by family.}
	\label{tab:family-summary}
	\begin{tabular}{@{}lrr@{}}
		\toprule
		Family & Species & Specimens \\
		\midrule
		Anacardiaceae & 2 & 10 \\
		Combretaceae  & 2 & 10 \\
		Fabaceae      & 3 & 15 \\
		Irvingiaceae  & 1 &  5 \\
		Malvaceae     & 3 & 15 \\
		Meliaceae     & 3 & 15 \\
		Moraceae      & 2 & 10 \\
		Sapotaceae    & 2 & 10 \\
		\midrule
		Total         & 18 & 90 \\
		\bottomrule
	\end{tabular}
\end{table}

\begin{table}[htbp]
	\centering
	\caption{Specimens by collector and phenological state.}
	\label{tab:collector-summary}
	\begin{tabular}{@{}llrrrr@{}}
		\toprule
		Collector & Herbarium & Flowering & Fruiting & Sterile & Total \\
		\midrule
		Harding   & \textsc{okf} &  8 & 19 &  5 & 32 \\
		Olaniyan  & \textsc{svn} & 11 & 16 &  5 & 32 \\
		Alade     & \textsc{auh} &  4 & 12 & 10 & 26 \\
		\midrule
		Total     &              & 23 & 47 & 20 & 90 \\
		\bottomrule
	\end{tabular}
\end{table}